\documentclass[11pt]{article}

\usepackage[margin=0.8in]{geometry}
\usepackage{amsmath, amssymb}
\usepackage{xcolor}
\usepackage{booktabs}
\usepackage{tabularx}
\usepackage{array}
\usepackage{tikz}
\usepackage{needspace}

\definecolor{newc}{HTML}{2E7D32}
\definecolor{tgtc}{HTML}{C62828}

\setlength{\parindent}{0pt}
\setlength{\parskip}{0.35em plus 0.1em minus 0.1em}
\newcolumntype{Y}{>{\raggedright\arraybackslash}X}

\tikzset{bn/.style={circle, draw, thick, fill=white, minimum size=5.6mm, inner sep=0pt, font=\footnotesize}}

% \cbt[layout levels]{values in level order}{indices with green ring}{indices with red ring}
% Node i is drawn at level floor(log2(i+1)); its parent is index (i-1)/2.
\newcommand{\cbt}[4][3]{%
  \begin{tikzpicture}[x=1mm, y=1mm, baseline=(current bounding box.north)]
    \ifcase#1\or\def\cbtW{16}\or\def\cbtW{26}\or\def\cbtW{34}\or\def\cbtW{66}\fi
    \def\cbtN{0}%
    \if\relax\detokenize{#2}\relax\else\foreach \cbtV [count=\cbtC] in {#2} {\xdef\cbtN{\cbtC}}\fi
    \pgfmathsetmacro{\cbtH}{9*((\cbtN>1)+(\cbtN>3)+(\cbtN>7))}
    \useasboundingbox (0,5) rectangle (\cbtW,-\cbtH-5);
    \if\relax\detokenize{#2}\relax
      \node[font=\footnotesize\itshape] at (\cbtW/2,-\cbtH/2) {(empty)};
    \else
      \foreach \cbtV [count=\cbtK from 0] in {#2} {
        \pgfmathtruncatemacro{\cbtL}{(\cbtK>0)+(\cbtK>2)+(\cbtK>6)+(\cbtK>14)}
        \pgfmathsetmacro{\cbtX}{(\cbtK-(2^\cbtL-1)+0.5)*\cbtW/(2^\cbtL)}
        \node[bn] (n\cbtK) at (\cbtX,-9*\cbtL) {\cbtV};
        \ifnum\cbtK>0
          \pgfmathtruncatemacro{\cbtP}{(\cbtK-1)/2}
          \draw[thick] (n\cbtP) -- (n\cbtK);
        \fi
      }
    \fi
    \if\relax\detokenize{#3}\relax\else
      \foreach \cbtG in {#3} \draw[newc, line width=1.4pt] (n\cbtG.center) circle (3.9);
    \fi
    \if\relax\detokenize{#4}\relax\else
      \foreach \cbtG in {#4} \draw[tgtc, line width=1.4pt] (n\cbtG.center) circle (3.9);
    \fi
  \end{tikzpicture}}

\newcommand{\tcell}[2]{\begin{tabular}[t]{@{}c@{}}#2\\\footnotesize #1\end{tabular}}
\newcommand{\trow}[1]{\par\smallskip\noindent\null\hfill#1\hfill\null\par\smallskip}
\newcommand{\gring}{\tikz[baseline=-0.6ex]\draw[newc, line width=1.4pt] (0,0) circle (1.1mm);}
\newcommand{\rring}{\tikz[baseline=-0.6ex]\draw[tgtc, line width=1.4pt] (0,0) circle (1.1mm);}
\newcommand{\prob}[1]{\par\medskip\needspace{40mm}{\large\bfseries #1}\par\smallskip}
\newcommand{\pa}[2][0mm]{\par\smallskip\needspace{#1}\textbf{(#2)}\ }
\newcommand{\cbox}[2]{\begin{minipage}[t]{0.485\textwidth}\raggedright #1\par\smallskip\centering#2\end{minipage}}
\newcommand{\crow}[2]{\par\smallskip\needspace{40mm}\noindent#1\hfill#2\par\medskip}

\begin{document}

%=============================================================================
\prob{Problem 1: Creating the Tree}

$\mathit{arr} = \{5, 12, 8, 21, 3, 17, 9\}$, $n = 7$.
\hfill \gring\ node added / value moved \qquad \rring\ deletion target

\pa{a}
\begin{center}\small
\begin{tabular}{@{}ccccccl@{}}
\toprule
Step & Element & Index $i$ & Parent index $(i-1)/2$ & Parent value & Left / Right & Rule(s)\\
\midrule
1 & 5  & 0 & ---           & --- & root             & C1\\
2 & 12 & 1 & $(1-1)/2 = 0$ & 5   & Left ($i$ odd)   & C2, C3\\
3 & 8  & 2 & $(2-1)/2 = 0$ & 5   & Right ($i$ even) & C2, C4\\
4 & 21 & 3 & $(3-1)/2 = 1$ & 12  & Left ($i$ odd)   & C2, C3\\
5 & 3  & 4 & $(4-1)/2 = 1$ & 12  & Right ($i$ even) & C2, C4\\
6 & 17 & 5 & $(5-1)/2 = 2$ & 8   & Left ($i$ odd)   & C2, C3\\
7 & 9  & 6 & $(6-1)/2 = 2$ & 8   & Right ($i$ even) & C2, C4\\
\bottomrule
\end{tabular}
\end{center}

\trow{\tcell{1: add 5}{\cbt{5}{0}{}}\hfill
      \tcell{2: add 12}{\cbt{5,12}{1}{}}\hfill
      \tcell{3: add 8}{\cbt{5,12,8}{2}{}}\hfill
      \tcell{4: add 21}{\cbt{5,12,8,21}{3}{}}}
\trow{\tcell{5: add 3}{\cbt{5,12,8,21,3}{4}{}}\hfill
      \tcell{6: add 17}{\cbt{5,12,8,21,3,17}{5}{}}\hfill
      \tcell{7: add 9 (final)}{\cbt{5,12,8,21,3,17,9}{6}{}}}

\textbf{Final tree (level order):} 5 12 8 21 3 17 9

\pa{b}
21, 3, 17, 9 (indices 3--6) have NULL children. Rule \textbf{C5}: their child indices $2i+1$, $2i+2$
are 7--14, all $\ge n = 7$, so those children do not exist.

\par\smallskip
\begin{minipage}[t]{0.40\textwidth}
\textbf{(c)}\par\smallskip
\small
\begin{tabular}{@{}ll@{}}
\toprule
Root & 5\\
Leaves & 21, 3, 17, 9\\
Internal nodes & 5, 12, 8\\
Edges & 6 \ ($n - 1$)\\
Height & 2 \ ($5 \to 12 \to 21$)\\
Depth of 17 & 2\\
Ancestors of 3 & 12, 5\\
Descendants of 12 & 21, 3\\
Subtree at 8 & $\{8, 17, 9\}$\\
Path 5 to 9 & $5 \to 8 \to 9$, length 2\\
\bottomrule
\end{tabular}
\end{minipage}\hfill
\begin{minipage}[t]{0.56\textwidth}
\textbf{(d)}\par\smallskip
\small
\begin{tabularx}{\linewidth}{@{}lY@{}}
\toprule
Max nodes on level 2 & $2^2 = 4$; level 2 has 4 \checkmark\\
Max nodes for $H = 2$ & $2^{3} - 1 = 7 = n$ \checkmark\\
Leaves $=$ two-child $+\,1$ & 5, 12, 8 $\Rightarrow 3 + 1 = 4$ \checkmark\\
Min height, $N = 7$ & $\lfloor \log_2 7 \rfloor = 2 =$ height \checkmark\\
\bottomrule
\end{tabularx}

\normalsize\medskip
\textbf{Both full and complete.} C6: creation fills level by level, left to right, so the tree is
always complete. Since $n = 7 = 2^{H+1} - 1$, every level is filled, so it is also full.
\end{minipage}

%=============================================================================
\prob{Problem 2: Traversing the Tree}

\pa{a--c} Each cell gives the node visited and why it is visited at that moment.
\begin{center}\small
\begin{tabularx}{\textwidth}{@{}cYYY@{}}
\toprule
\# & (a) Preorder (N L R) & (b) Inorder (L N R) & (c) Postorder (L R N)\\
\midrule
1 & \textbf{5}: root; N, then L (12), R (8) & \textbf{21}: L is NULL, so N & \textbf{21}: L, R NULL, so N\\
2 & \textbf{12}: N, then L (21), R (3)       & \textbf{12}: L (21) done, so N & \textbf{3}: L, R NULL, so N\\
3 & \textbf{21}: N; L, R NULL                & \textbf{3}: L is NULL, so N & \textbf{12}: L (21), R (3) done, so N\\
4 & \textbf{3}: N; L, R NULL                 & \textbf{5}: L (12) done, so N & \textbf{17}: L, R NULL, so N\\
5 & \textbf{8}: N, then L (17), R (9)        & \textbf{17}: L is NULL, so N & \textbf{9}: L, R NULL, so N\\
6 & \textbf{17}: N; L, R NULL                & \textbf{8}: L (17) done, so N & \textbf{8}: L (17), R (9) done, so N\\
7 & \textbf{9}: N; L, R NULL                 & \textbf{9}: L is NULL, so N & \textbf{5}: L (12), R (8) done, so N\\
\midrule
Rules & T2; T1, T6 at leaves & T3; T1, T6 at leaves & T4; T1, T6 at leaves\\
Result & \textbf{5 12 21 3 8 17 9} & \textbf{21 12 3 5 17 8 9} & \textbf{21 3 12 17 9 8 5}\\
\bottomrule
\end{tabularx}
\end{center}

\pa[45mm]{d} Level order (T5 at every row).
\begin{center}\small
\begin{tabular}{@{}clclll@{}}
\toprule
\# & Queue before & Remove \& visit & Add children (left, right) & Queue after & Output so far\\
\midrule
1 & 5            & 5  & 12, 8 & 12, 8        & 5\\
2 & 12, 8        & 12 & 21, 3 & 8, 21, 3     & 5 12\\
3 & 8, 21, 3     & 8  & 17, 9 & 21, 3, 17, 9 & 5 12 8\\
4 & 21, 3, 17, 9 & 21 & none  & 3, 17, 9     & 5 12 8 21\\
5 & 3, 17, 9     & 3  & none  & 17, 9        & 5 12 8 21 3\\
6 & 17, 9        & 17 & none  & 9            & 5 12 8 21 3 17\\
7 & 9            & 9  & none  & empty, stop  & \textbf{5 12 8 21 3 17 9}\\
\bottomrule
\end{tabular}
\end{center}

\pa{e}
The tree was built in level order (C6), so T5 visits the nodes in index order $0, 1, \dots, 6$,
which is \textit{arr}. \textbf{Postorder} visits the root last: by T4 a node is visited only after
both of its subtrees, and every other node is in a subtree of 5.

\pa{f}
Search 3: compares 5, 12, 8, 21, 3, so it is found after \textbf{5} comparisons.
Search 4: compares all 7 nodes, so it is not found after \textbf{7} comparisons (worst case, $n$).

%=============================================================================
\prob{Problem 3: Adding Elements to the Tree}

Each insertion after (a) scans in level order with a queue (I2). Children are listed as
(left, right); -- means NULL.
\begin{center}\small
\begin{tabularx}{\textwidth}{@{}cYlccl@{}}
\toprule
Insert & Nodes checked (left, right) & Rule per node & Parent & Side & Level order after\\
\midrule
11 & 5 (12, 8), 12 (21, 3), 8 (17, 9), 21 (--, --) & I5 $\times$3, I3 & 21 & Left & 5 12 8 21 3 17 9 11\\
2  & 5, 12, 8 (both present), 21 (11, --) & I5 $\times$3, I4 & 21 & Right & 5 12 8 21 3 17 9 11 2\\
14 & 5, 12, 8, 21 (both present), 3 (--, --) & I5 $\times$4, I3 & 3 & Left & 5 12 8 21 3 17 9 11 2 14\\
6  & 5, 12, 8, 21 (both present), 3 (14, --) & I5 $\times$4, I4 & 3 & Right & 5 12 8 21 3 17 9 11 2 14 6\\
\bottomrule
\end{tabularx}
\end{center}

\textbf{(a)} The tree is empty, so 5 becomes the root (\textbf{I1}).\\
\textbf{(b)} 5, 12, 8 are full (I5); 21 has a free left position, so 11 is the \textbf{left child of 21} (I3).\\
\textbf{(c)} 21's left is now 11, so I3 fails and \textbf{I4} puts 2 on the \textbf{right of 21}
(index $8 = 2(3) + 2$).\\
\textbf{(d)} Skipped by I5: \textbf{5, 12, 8, 21}. 14 is the \textbf{left child of 3} (I3).\\
\textbf{(e)} I6: index $= n = 10$; parent $(10-1)/2 = 4$, value 3; $10 = 2(4)+2$, so \textbf{right of 3}.
Full scan: 3 has left 14 and a free right (I4) \checkmark

\trow{\tcell{(a) insert 5}{\cbt[1]{5}{0}{}}\hfill
      \tcell{(b) insert 11}{\cbt[4]{5,12,8,21,3,17,9,11}{7}{}}\hfill
      \tcell{(c) insert 2}{\cbt[4]{5,12,8,21,3,17,9,11,2}{8}{}}}
\trow{\tcell{(d) insert 14}{\cbt[4]{5,12,8,21,3,17,9,11,2,14}{9}{}}\hfill
      \tcell{(e) insert 6}{\cbt[4]{5,12,8,21,3,17,9,11,2,14,6}{10}{}}}

%=============================================================================
\prob{Problem 4: Deleting Elements (Every Case)}

Each part starts from the original tree 5 12 8 21 3 17 9, whose deepest node is 9 (D3).

\crow{\cbox{\textbf{(a) Case A: delete 4.} Level order checks all 7 nodes and never finds 4 (D3),
so the tree is unchanged (D4).}
{\tcell{Before = after}{\cbt{5,12,8,21,3,17,9}{}{}}}}
{\cbox{\textbf{(b) Case B: delete 9.} Target $=$ deepest (D7): nothing is copied; D6 clears 8's
right link. Easiest case: no copy, no other node changes.}
{\tcell{Before}{\cbt{5,12,8,21,3,17,9}{}{6}}\hspace{6mm}\tcell{After}{\cbt{5,12,8,21,3,17}{}{}}}}

\crow{\cbox{\textbf{(c) Case C: delete 21.} Target 21 (leaf), deepest 9. Copy 9 into 21's node (D5),
remove the old 9 (D6).}
{\tcell{Before}{\cbt{5,12,8,21,3,17,9}{}{3}}\hspace{6mm}\tcell{After}{\cbt{5,12,8,9,3,17}{3}{}}}}
{\cbox{\textbf{(d) Case D: delete 12.} Target 12 (internal), deepest 9 (D5, D6). The replacement 9
keeps 12's children: \textbf{21 (left) and 3 (right)}.}
{\tcell{Before}{\cbt{5,12,8,21,3,17,9}{}{1}}\hspace{6mm}\tcell{After}{\cbt{5,9,8,21,3,17}{1}{}}}}

\crow{\cbox{\textbf{(e) Case E: delete 5.} Target 5 (root), deepest 9. Copy 9 into the root (D5),
remove the old 9 (D6). New root: \textbf{9}.}
{\tcell{Before}{\cbt{5,12,8,21,3,17,9}{}{0}}\hspace{6mm}\tcell{After}{\cbt{9,12,8,21,3,17}{0}{}}}}
{\cbox{\textbf{(f) Case F: $\{9\}$.} Delete 9: the only node holds the key, so the tree becomes empty.
Delete 4 on a fresh $\{9\}$: not found, so it stays $\{9\}$ (D2).\par
\textbf{(g) Case G:} the tree is empty; delete 5 does nothing (D1).}
{\tcell{$\{9\}$}{\cbt[1]{9}{}{0}}\hspace{1mm}\tcell{del 9}{\cbt[1]{}{}{}}\hspace{3mm}
 \tcell{$\{9\}$}{\cbt[1]{9}{}{}}\hspace{1mm}\tcell{del 4}{\cbt[1]{9}{}{}}}}

\pa[40mm]{h}
Removing the target directly would leave a hole: an internal target's children would lose their
parent, and the array would have a gap, so the tree would no longer be complete. The deepest node is
the last node in level order and always a leaf, so removing it disconnects nothing and keeps the tree
complete; only one value is copied (D5) and one link cleared (D6).

\begin{center}\small
\begin{tabular}{@{}cccccl@{}}
\toprule
Case & Delete & Target & Deepest & Rules applied & Level order after\\
\midrule
A & 4     & not found     & 9     & D3, D4     & 5 12 8 21 3 17 9 (unchanged)\\
B & 9     & 9             & 9     & D3, D7     & 5 12 8 21 3 17\\
C & 21    & 21            & 9     & D3, D5, D6 & 5 12 8 9 3 17\\
D & 12    & 12            & 9     & D3, D5, D6 & 5 9 8 21 3 17\\
E & 5     & 5             & 9     & D3, D5, D6 & 9 12 8 21 3 17\\
F & 9 / 4 & 9 / not found & 9 / 9 & D2         & empty / 9\\
G & 5     & ---           & ---   & D1         & empty\\
\bottomrule
\end{tabular}
\end{center}

%=============================================================================
\prob{Problem 5: Deleting All Elements}

\pa{a} Delete 5, 12, 8, 21, 3, 17, 9 without resetting the tree.
\begin{center}\small
\begin{tabular}{@{}cclcccl@{}}
\toprule
\# & Delete & Tree before & Deepest & Case & Rules & Tree after\\
\midrule
1 & 5  & 5 12 8 21 3 17 9 & 9  & E (root)     & D3, D5, D6 & 9 12 8 21 3 17\\
2 & 12 & 9 12 8 21 3 17   & 17 & D (internal) & D3, D5, D6 & 9 17 8 21 3\\
3 & 8  & 9 17 8 21 3      & 3  & C (leaf)     & D3, D5, D6 & 9 17 3 21\\
4 & 21 & 9 17 3 21        & 21 & B (deepest)  & D3, D7     & 9 17 3\\
5 & 3  & 9 17 3           & 3  & B (deepest)  & D3, D7     & 9 17\\
6 & 17 & 9 17             & 17 & B (deepest)  & D3, D7     & 9\\
7 & 9  & 9                & 9  & F (one node) & D2         & empty\\
\bottomrule
\end{tabular}
\end{center}

\pa{b}
\trow{\tcell{Start}{\cbt{5,12,8,21,3,17,9}{}{}}\hfill
      \tcell{1: del 5}{\cbt{9,12,8,21,3,17}{0}{}}\hfill
      \tcell{2: del 12}{\cbt{9,17,8,21,3}{1}{}}\hfill
      \tcell{3: del 8}{\cbt{9,17,3,21}{2}{}}}
\trow{\tcell{4: del 21}{\cbt{9,17,3}{}{}}\hfill
      \tcell{5: del 3}{\cbt{9,17}{}{}}\hfill
      \tcell{6: del 17}{\cbt{9}{}{}}\hfill
      \tcell{7: del 9}{\cbt[1]{}{}{}}}

\pa{c}
Cases E, D, C, B (three times) and F appeared; A and G did not (every key was present and the tree was
never empty at the start of a deletion). The tree stayed complete after every step, because each
deletion only removes the last node in level order, so the remaining nodes fill indices 0 to $n-2$
with no gaps. Values copied: \textbf{3}.

%=============================================================================
\prob{Bonus: Reverse Order}

Delete 9, 17, 3, 21, 8, 12, 5 from the original tree.
\begin{center}\small
\begin{tabular}{@{}cclcccl@{}}
\toprule
\# & Delete & Tree before & Deepest & Case & Rules & Tree after\\
\midrule
1 & 9  & 5 12 8 21 3 17 9 & 9  & B & D3, D7 & 5 12 8 21 3 17\\
2 & 17 & 5 12 8 21 3 17   & 17 & B & D3, D7 & 5 12 8 21 3\\
3 & 3  & 5 12 8 21 3      & 3  & B & D3, D7 & 5 12 8 21\\
4 & 21 & 5 12 8 21        & 21 & B & D3, D7 & 5 12 8\\
5 & 8  & 5 12 8           & 8  & B & D3, D7 & 5 12\\
6 & 12 & 5 12             & 12 & B & D3, D7 & 5\\
7 & 5  & 5                & 5  & F & D2     & empty\\
\bottomrule
\end{tabular}
\end{center}

\trow{\tcell{1: del 9}{\cbt{5,12,8,21,3,17}{}{}}\hfill
      \tcell{2: del 17}{\cbt{5,12,8,21,3}{}{}}\hfill
      \tcell{3: del 3}{\cbt{5,12,8,21}{}{}}\hfill
      \tcell{4: del 21}{\cbt{5,12,8}{}{}}}
\trow{\tcell{5: del 8}{\cbt{5,12}{}{}}\hfill
      \tcell{6: del 12}{\cbt{5}{}{}}\hfill
      \tcell{7: del 5}{\cbt[1]{}{}{}}}

Case B at steps 1--6 and F at step 7. This order is the simplest because reversing \textit{arr}
reverses the level order (C6), so the key deleted is always the deepest node: nothing is copied and
only one leaf link is cleared (D7). Values copied (D5): \textbf{3} in the forward order, \textbf{0} in
the reverse order.

\end{document}
